% ---------------------------------------------------------------------------
% Rewritten in the Day 6 review. The previous version led with "Lemma 1 /
% Lemma 2 / Theorem 1" and the contraction-mapping vocabulary. The convergence
% story is genuinely just "Day 2 policy iteration, unchanged", so it is now
% told that way; the paper's numbering is kept only as a footnote pointer.
% ---------------------------------------------------------------------------
\begin{frame}{Soft Policy Iteration --- The Plan}
\begin{itemize}
    \item You already know \textbf{policy iteration} from Day 2: alternate two steps until nothing moves
    \begin{enumerate}
        \item \textbf{evaluation:} given $\pi$, work out its $Q$
        \item \textbf{improvement:} given $Q$, make $\pi$ greedier
    \end{enumerate}
    \vspace{0.6em}
    \item \textbf{Soft policy iteration is the same loop with the dictionary applied} --- ``greedier'' becomes ``closer to the softmax of $Q$'', and the values carry the entropy bonus
    \vspace{0.6em}
    \item The next two slides do one step each. The only question worth asking is: \emph{does swapping $\max$ for soft max break the guarantees from Day 2?}
    \item[$\Rightarrow$] \textcolor{red}{It does not.} And once we know that, we can go build the deep version (SAC) without worrying
\end{itemize}
\end{frame}

\begin{frame}{Step 1 --- Soft Policy Evaluation}
\begin{itemize}
    \item \textbf{Goal:} given a fixed $\pi$, compute its value \emph{under the MaxEnt objective}. Same Bellman backup as Day 2, with the entropy bonus folded into the value:
\end{itemize}
\begin{align*}
\mathcal{T}^{\pi} Q(\mathbf{s}_t, \mathbf{a}_t) &= r(\mathbf{s}_t, \mathbf{a}_t) + \gamma\, \mathbb{E}_{\mathbf{s}_{t+1} \sim p}\!\left[\, V(\mathbf{s}_{t+1})\, \right] \\[0.3em]
V(\mathbf{s}_t) &= \mathbb{E}_{\mathbf{a}_t \sim \pi}\!\big[\, Q(\mathbf{s}_t, \mathbf{a}_t) \; \textcolor{red}{-\; \alpha \log \pi(\mathbf{a}_t \mid \mathbf{s}_t)} \,\big]
\;=\; \mathbb{E}_{\pi}[Q] + \alpha\,\mathcal{H}(\pi(\cdot\mid\mathbf{s}_t))
\end{align*}
\begin{itemize}
    \item \textbf{Does repeating this converge?} Yes --- and for a boring reason. Compare with Day 2:
    \begin{itemize}
        \item Day 2's proof only needed the backup to \emph{shrink differences} between value estimates by a factor $\gamma$ each sweep. That is what makes repeated application settle on one answer
        \item The single new term $-\alpha\log\pi$ does not depend on $Q$ at all. It is just \textbf{an extra piece of reward}, and the Day 2 proof never cared what the reward was
    \end{itemize}
    \item[$\Rightarrow$] \textbf{Nothing new to prove.} Soft policy evaluation is ordinary policy evaluation on the reward $r + \alpha\mathcal{H}$
\end{itemize}
\end{frame}

\begin{frame}{Step 2 --- Soft Policy Improvement}
\small
\begin{itemize}
    \item \textbf{Day 2:} improve $\pi$ by making it greedy, $\pi \leftarrow \arg\max_a Q(s,a)$
    \item \textbf{Today:} improve $\pi$ by moving it toward the \emph{softmax} of $Q$, $\;\pi \propto \exp(\tfrac1\alpha Q)$ --- the distribution we showed is optimal for ``value $+\,\alpha\,$entropy''
    \vspace{0.4em}
    \item \textcolor{red}{One practical snag:} $\exp(\tfrac1\alpha Q)$ can come out \emph{any shape at all}, and no network can output an arbitrary distribution. Whatever actor we build will only produce distributions from some limited family --- call it $\Pi$
    \item So we settle for \emph{the closest member of $\Pi$}, and ``closest'' between two distributions means the KL divergence you met in TRPO (Day 4):
\end{itemize}
\[
\pi_{\text{new}} = \arg\min_{\pi' \in \Pi} \; D_{\mathrm{KL}}\!\Big( \pi'(\,\cdot \mid \mathbf{s}_t) \,\Big\|\, \underbrace{\exp\!\big( \tfrac{1}{\alpha} Q^{\pi_{\text{old}}}(\mathbf{s}_t, \cdot)\big) \big/ Z(\mathbf{s}_t)}_{\text{the softmax policy we are aiming at}} \Big)
\]
\vspace{-0.3em}
\begin{itemize}
    \item[] {\footnotesize $\Pi$ stays abstract for now --- the theory works for \emph{any} choice; we pick a concrete one when we build SAC. $Z(\mathbf{s}_t)$ just makes the target sum to one, and drops out of the $\arg\min$ because it does not depend on $\pi'$.}
    \vspace{0.4em}
    \item \textbf{Does this actually improve anything?} Yes, and again for a simple reason: \emph{$\pi_{\text{old}}$ is itself one of the candidates} in the $\arg\min$. So $\pi_{\text{new}}$ is by construction no worse --- the soft value can only go up
\end{itemize}
\end{frame}

\begin{frame}{Putting It Together}
\begin{itemize}
    \item \textbf{Soft policy iteration} $=$ alternate the two steps: evaluate $\longleftrightarrow$ improve
    \vspace{0.4em}
    \item \textbf{Why it lands somewhere:} improvement means the values only ever go \emph{up}; rewards and entropy are bounded so they cannot go up forever. Something that only climbs and cannot climb past a ceiling has to settle --- and where it settles is the best MaxEnt policy the network family can express
    \vspace{0.6em}
    \item \textcolor{red}{But we cannot run this loop as written.} It assumes we can sweep over \emph{every} state and action --- fine for a Day 2 gridworld, hopeless for a robot with continuous joints
    \vspace{0.4em}
    \item[$\Rightarrow$] Replace the two exact steps with two \textbf{neural networks trained by gradient descent}:
    \begin{itemize}
        \item evaluation $\;\to\;$ a critic $Q_\phi$ fitting the soft Bellman target
        \item improvement $\;\to\;$ an actor $\pi_\theta$ taking gradient steps on that KL
    \end{itemize}
    \item[] \textcolor{red}{That is Soft Actor-Critic.}
\end{itemize}
\vspace{0.1em}
{\footnotesize The two ``does it work'' arguments above are Lemmas 1--2 in the paper \rref{1}, and the last one is its Theorem 1 --- look there if you want them stated formally.}
\end{frame}
